\chapter{2026年全国II卷}

\section{单选题}

\begin{problem}
\(\left( 1-3\i \right)^2=\)（\quad）

\choices
    {\(-8+6\i\)}
    {\(-8-6\i\)}
    {\(8+6\i\)}
    {\(8-6\i\)}

\end{problem}

\begin{answer}
B
\end{answer}

\begin{solution}
因为 \(\left( 1-3\i \right)^2=1-6\i+9\i^2=1-6\i-9=-8-6\i\)，所以选：B.
\end{solution}

\begin{problem}
已知集合 \(A=\{0,1,3,6,9\}\)，\(B=\{x \mid \sqrt{x}=x\}\)，则 \(A\cap B=\)（\quad）

\choices
    {\(\{0,1\}\)}
    {\(\{3,6\}\)}
    {\(\{0,1,9\}\)}
    {\(\{0,3,9\}\)}

\end{problem}

\begin{answer}
A
\end{answer}

\begin{solution}
由 \(\sqrt{x}=x\) 可知 \(x\ge 0\). 两边平方，得 \(x=x^2\)，即 \(x\left( x-1 \right)=0\). 所以 \(x=0\) 或 \(x=1\)，从而 \(B=\{0,1\}\). 因此 \(A\cap B=\{0,1\}\). 故选：A.
\end{solution}

\begin{problem}
已知向量 \(\bs{a}\)，\(\bs{b}\) 满足 \(\left| \bs{a}+\bs{b} \right|=1\)，\(\left| \bs{a}-\bs{b} \right|=\sqrt{3}\)，则 \(\bs{a}\cdot \bs{b}=\)（\quad）

\choices
    {\(\frac{1}{2}\)}
    {\(\frac{1}{3}\)}
    {\(-\frac{1}{3}\)}
    {\(-\frac{1}{2}\)}

\end{problem}

\begin{answer}
D
\end{answer}

\begin{solution}
由向量数量积公式，得 \(\left| \bs{a}+\bs{b} \right|^2=\lvert \bs{a} \rvert^2+\lvert \bs{b} \rvert^2+2\bs{a}\cdot \bs{b}=1\)，\(\left| \bs{a}-\bs{b} \right|^2=\lvert \bs{a} \rvert^2+\lvert \bs{b} \rvert^2-2\bs{a}\cdot \bs{b}=3\). 两式相减，得 \(4\bs{a}\cdot \bs{b}=1-3=-2\). 所以 \(\bs{a}\cdot \bs{b}=-\frac{1}{2}\). 故选：D.
\end{solution}

\begin{problem}
设双曲线 \(C:\frac{x^2}{a^2}-\frac{y^2}{b^2}=1\)（\(a>0\)，\(b>0\)）经过点 \(\left( 1,0 \right)\) 和点 \(\left( \frac{\sqrt{7}}{2},3 \right)\)，则 \(C\) 的渐近线方程为（\quad）

\choices
    {\(y=\pm 3\sqrt{2}x\)}
    {\(y=\pm 2\sqrt{3}x\)}
    {\(y=\pm \frac{\sqrt{3}}{6}x\)}
    {\(y=\pm \frac{\sqrt{2}}{6}x\)}

\end{problem}

\begin{answer}
B
\end{answer}

\begin{solution}
因为点 \(\left( 1,0 \right)\) 在双曲线上，所以 \(\frac{1}{a^2}=1\)，故 \(a^2=1\).
再由点 \(\left( \frac{\sqrt{7}}{2},3 \right)\) 在双曲线上，得 \(\frac{\left( \frac{\sqrt{7}}{2} \right)^2}{1}-\frac{3^2}{b^2}=1\). 化简得 \(\frac{7}{4}-\frac{9}{b^2}=1\)，所以 \(\frac{9}{b^2}=\frac{3}{4}\)，从而 \(b^2=12\). 双曲线的渐近线方程为 \(y=\pm \frac{b}{a}x=\pm 2\sqrt{3}x\). 故选：B.
\end{solution}

\begin{problem}
已知棱台的上，下底面均是有一个内角为 \(60^\circ\) 的菱形，上，下底面的边长分别为 \(2\)，\(3\)，该棱台的高为 \(\sqrt{3}\)，则其体积为（\quad）

\choices
    {\(\frac{19}{12}\)}
    {\(\frac{19}{6}\)}
    {\(\frac{19}{4}\)}
    {\(\frac{19}{2}\)}

\end{problem}

\begin{answer}
D
\end{answer}

\begin{solution}
边长为 \(a\)，内角为 \(60^\circ\) 的菱形面积为 \(a^2\sin 60^\circ=\frac{\sqrt{3}}{2}a^2\). 所以上底面积 \(S_1=\frac{\sqrt{3}}{2}\times 2^2=2\sqrt{3}\)，下底面积 \(S_2=\frac{\sqrt{3}}{2}\times 3^2=\frac{9\sqrt{3}}{2}\). 又 \(\sqrt{S_1S_2}=\sqrt{2\sqrt{3}\cdot \frac{9\sqrt{3}}{2}}=3\sqrt{3}\). 由棱台体积公式 \(V=\frac{h}{3}\left( S_1+S_2+\sqrt{S_1S_2} \right)\)，得 \(V=\frac{\sqrt{3}}{3}\left( 2\sqrt{3}+\frac{9\sqrt{3}}{2}+3\sqrt{3} \right)=\frac{\sqrt{3}}{3}\cdot \frac{19\sqrt{3}}{2}\). 所以 \(V=\frac{19}{2}\).
故选：D.
\end{solution}

\begin{problem}
把甲，乙，丙，丁等 \(8\) 位员工分为 \(A\)，\(B\) 两个技术攻关组，要求每组 \(4\) 人，其中甲，乙同组，丙，丁不同组，则分组方案共有（\quad）

\choices
    {\(10\) 种}
    {\(12\) 种}
    {\(16\) 种}
    {\(24\) 种}

\end{problem}

\begin{answer}
C
\end{answer}

\begin{solution}
因为 \(A\)，\(B\) 两组有标号，所以先确定甲，乙所在的组，有 \(2\) 种选法.
固定甲，乙在其中一组后，为了使丙，丁不同组，这一组在丙，丁中必须恰好选入 \(1\) 人，有 \(2\) 种选法. 其余 \(1\) 人从另外 \(4\) 人中任取 \(1\) 人，有 \(4\) 种选法.
因此分组方案总数为 \(2\times 2\times 4=16\). 故选：C.
\end{solution}

\begin{problem}
已知 \(\alpha\) 为第二象限角，\(3\sin 2\alpha \cos \alpha=8\sin \alpha \cos 2\alpha\)，则 \(\frac{1+\sin \alpha}{2-\cos \alpha}=\)（\quad）

\choices
    {\(\frac{3}{4}\)}
    {\(\frac{3}{5}\)}
    {\(\frac{1}{2}\)}
    {\(\frac{5}{12}\)}

\end{problem}

\begin{answer}
C
\end{answer}

\begin{solution}
由 \(\sin 2\alpha=2\sin \alpha \cos \alpha\)，原式化为 \(6\sin \alpha \cos^2 \alpha=8\sin \alpha \cos 2\alpha\). 因为 \(\alpha\) 为第二象限角，所以 \(\sin \alpha>0\). 两边同除以 \(2\sin \alpha\)，得 \(3\cos^2 \alpha=4\cos 2\alpha\). 又 \(\cos 2\alpha=2\cos^2 \alpha-1\)，所以 \(3\cos^2 \alpha=8\cos^2 \alpha-4\). 从而 \(5\cos^2 \alpha=4\)，即 \(\cos^2 \alpha=\frac{4}{5}\). 因为 \(\alpha\) 为第二象限角，所以 \(\cos \alpha=-\frac{2}{\sqrt{5}}\)，且 \(\sin \alpha=\frac{1}{\sqrt{5}}\). 于是 \(\frac{1+\sin \alpha}{2-\cos \alpha}=\frac{1+\frac{1}{\sqrt{5}}}{2+\frac{2}{\sqrt{5}}}=\frac{\sqrt{5}+1}{2\sqrt{5}+2}\). 进一步化简得 \(\frac{1+\sin \alpha}{2-\cos \alpha}=\frac{1}{2}\).
故选：C.
\end{solution}

\begin{problem}
已知 \(f\left( x \right)\) 是定义域为 \(\R\) 的偶函数，\(f\left( x \right)+f\left( x-2 \right)=0\)，且当 \(x\in \left[ \frac{3}{2},3 \right]\) 时，\(f\left( x \right)=x^2+ax+b\)，则

\choices
    {\(a=-2\)，\(b=-3\)}
    {\(a=-2\)，\(b=3\)}
    {\(a=-4\)，\(b=-3\)}
    {\(a=-4\)，\(b=3\)}

\end{problem}

\begin{answer}
D
\end{answer}

\begin{solution}
设 \(p(x)=x^2+ax+b\). 当 \(x\in \left[ \frac{3}{2},3 \right]\) 时，\(f(x)=p(x)\).
由 \(f(x)+f(x-2)=0\)，取 \(x=1\)，得 \(f(1)+f(-1)=0\). 又 \(f(x)\) 为偶函数，所以 \(f(-1)=f(1)\). 因此 \(2f(1)=0\)，即 \(f(1)=0\). 再取 \(x=3\)，得 \(f(3)+f(1)=0\)，所以 \(f(3)=0\). 因为 \(3\in \left[ \frac{3}{2},3 \right]\)，故 \(p(3)=0\)，即 \(9+3a+b=0\).
再取 \(x=\frac{3}{2}\)，得 \(f\left( \frac{3}{2} \right)+f\left( -\frac{1}{2} \right)=0\). 由偶函数性质，\(f\left( -\frac{1}{2} \right)=f\left( \frac{1}{2} \right)\)，所以 \(f\left( \frac{1}{2} \right)=-p\left( \frac{3}{2} \right)\). 再取 \(x=\frac{5}{2}\)，得 \(f\left( \frac{5}{2} \right)+f\left( \frac{1}{2} \right)=0\). 于是 \(p\left( \frac{5}{2} \right)=p\left( \frac{3}{2} \right)\)，即 \(\left( \frac{5}{2} \right)^2+\frac{5}{2}a+b=\left( \frac{3}{2} \right)^2+\frac{3}{2}a+b\).
化简得 \(4+a=0\)，所以 \(a=-4\). 代入 \(9+3a+b=0\)，得 \(b=3\). 故选：D.
\end{solution}
\section{多选题}

\begin{problem}
已知 \(\odot O:x^2+y^2=1\)，\(\odot A:x^2+y^2-6x-8y+k=0\)，则（\quad）

\choices
    {点 \(A\) 的坐标为 \(\left( -3,-4 \right)\)}
    {当 \(k=9\) 时，\(\odot A\) 与 \(x\) 轴相切}
    {当 \(k=-11\) 时，\(\odot O\) 与 \(\odot A\) 相切}
    {当 \(\odot O\) 与 \(\odot A\) 相交时，两个交点所在直线的方程为 \(6x+8y-k-2=0\)}

\end{problem}

\begin{answer}
BC
\end{answer}

\begin{solution}
将 \(\odot A\) 的方程配方，得 \(\left( x-3 \right)^2+\left( y-4 \right)^2=25-k\). 所以圆心 \(A\) 的坐标为 \(\left( 3,4 \right)\)，故 A 错误.

当 \(k=9\) 时，\(\odot A\) 的半径为 \(r=\sqrt{25-9}=4\). 圆心 \(A\left( 3,4 \right)\) 到 \(x\) 轴的距离也是 \(4\)，所以 \(\odot A\) 与 \(x\) 轴相切，B 正确.

当 \(k=-11\) 时，\(\odot A\) 的半径为 \(r=\sqrt{25-(-11)}=6\). 圆 \(\odot O\) 的半径为 \(1\)，两圆圆心距为 \(\sqrt{3^2+4^2}=5\). 又 \(6-1=5\)，所以两圆内切，C 正确.

若两圆相交，则两圆交点满足 \(x^2+y^2=1\) 和 \(x^2+y^2-6x-8y+k=0\). 两式相减，得 \(6x+8y-k-1=0\). 因此 D 中的直线方程不正确，D 错误.
故选：BC.
\end{solution}

\begin{problem}
设等比数列 \(\{a_n\}\) 的公比 \(q\ne 1\)，\(a_1>0\)，\(2a_3=a_1+a_2\)，记其前 \(n\) 项和为 \(S_n\)，则（\quad）

\choices
    {\(q=-\frac{1}{2}\)}
    {\(S_n>\frac{2a_1}{3}\)}
    {\(2S_{n+2}=S_{n+1}+S_n\)}
    {\(\sum_{k=1}^{n} S_k>\frac{2na_1}{3}\)}

\end{problem}

\begin{answer}
ACD
\end{answer}

\begin{solution}
因为 \(\{a_n\}\) 是等比数列，所以 \(a_2=a_1q\)，\(a_3=a_1q^2\). 由 \(2a_3=a_1+a_2\)，得 \(2a_1q^2=a_1+a_1q\). 又 \(a_1>0\)，所以 \(2q^2-q-1=0\). 解得 \(q=1\) 或 \(q=-\frac{1}{2}\). 因为 \(q\ne 1\)，故 \(q=-\frac{1}{2}\). 所以 A 正确.
此时 \(S_n=\frac{a_1\left( 1-q^n \right)}{1-q}=\frac{2a_1}{3}\left( 1-\left( -\frac{1}{2} \right)^n \right)\).
对于 B，取 \(n=2\)，得 \(S_2=\frac{2a_1}{3}\left( 1-\frac{1}{4} \right)=\frac{a_1}{2}<\frac{2a_1}{3}\). 所以 B 错误.
对于 C，\(2S_{n+2}=\frac{4a_1}{3}\left( 1-\left( -\frac{1}{2} \right)^{n+2} \right)\)，而 \(S_{n+1}+S_n=\frac{2a_1}{3}\left( 2-\left( -\frac{1}{2} \right)^{n+1}-\left( -\frac{1}{2} \right)^n \right)\). 又 \(\left( -\frac{1}{2} \right)^{n+1}+\left( -\frac{1}{2} \right)^n=\left( -\frac{1}{2} \right)^n\left( 1-\frac{1}{2} \right)=\frac{1}{2}\left( -\frac{1}{2} \right)^n\)，
并且 \(2\left( -\frac{1}{2} \right)^{n+2}=\frac{1}{2}\left( -\frac{1}{2} \right)^n\). 因此 \(2S_{n+2}=S_{n+1}+S_n\). 所以 C 正确.
对于 D，设 \(T_n=\sum_{k=1}^{n} S_k\). 则 \(T_n=\frac{2a_1}{3}\sum_{k=1}^{n}\left( 1-\left( -\frac{1}{2} \right)^k \right)=\frac{2na_1}{3}-\frac{2a_1}{3}\sum_{k=1}^{n}\left( -\frac{1}{2} \right)^k\).
又 \(\sum_{k=1}^{n}\left( -\frac{1}{2} \right)^k=-\frac{1}{3}\left( 1-\left( -\frac{1}{2} \right)^n \right)\)，所以 \(T_n=\frac{2na_1}{3}+\frac{2a_1}{9}\left( 1-\left( -\frac{1}{2} \right)^n \right)\). 因为 \(1-\left( -\frac{1}{2} \right)^n>0\)，故 \(T_n>\frac{2na_1}{3}\). 所以 D 正确.
故选：ACD.
\end{solution}

\begin{problem}
已知抛物线 \(E:y^2=8x\)，斜率为 \(k\)（\(k>0\)）的直线 \(l\) 经过点 \(\left( -1,0 \right)\)，等边三角形 \(ABC\) 的顶点 \(A\) 在 \(E\) 上，顶点 \(B\)，\(C\) 均在 \(l\) 上. 下列结论正确的有（\quad）

\choices
    {\(E\) 的准线方程为 \(x=-2\)}
    {若 \(l\) 与 \(E\) 没有公共点，则 \(k>\sqrt{2}\)}
    {若 \(l\) 与 \(E\) 的唯一公共点为 \(B\)，则 \(E\) 的焦点在直线 \(AB\) 上}
    {若 \(k=2\)，则 \(\triangle ABC\) 面积的最小值为 \(\frac{\sqrt{3}}{15}\)}

\end{problem}

\begin{answer}
ABD
\end{answer}

\begin{solution}
抛物线 \(E:y^2=8x\) 可写成 \(y^2=4px\) 的形式，其中 \(p=2\). 因此其准线方程为 \(x=-2\)，故 A 正确.
直线 \(l\) 的方程为 \(y=k\left( x+1 \right)\). 将其代入 \(y^2=8x\)，得 \(k^2\left( x+1 \right)^2=8x\)，即 \(k^2x^2+\left( 2k^2-8 \right)x+k^2=0\). 若 \(l\) 与 \(E\) 没有公共点，则上述方程无实根，所以判别式小于 \(0\)，即 \(\left( 2k^2-8 \right)^2-4k^4<0\). 化简得 \(64-32k^2<0\)，所以 \(k^2>2\). 又 \(k>0\)，故 \(k>\sqrt{2}\). 所以 B 正确.

若 \(l\) 与 \(E\) 的唯一公共点为 \(B\)，则上式判别式等于 \(0\)，从而 \(k=\sqrt{2}\). 此时交点为切点 \(B\left( 1,2\sqrt{2} \right)\)，抛物线的焦点为 \(F\left( 2,0 \right)\). 若焦点 \(F\) 在直线 \(AB\) 上，则直线 \(AB\) 的斜率应为 \(\frac{0-2\sqrt{2}}{2-1}=-2\sqrt{2}\). 而 \(BC\subset l\)，其斜率为 \(\sqrt{2}\). 由于 \(\triangle ABC\) 为等边三角形，\(\angle ABC=60^\circ\). 但若 \(AB\) 的斜率为 \(-2\sqrt{2}\)，则 \(\left| \frac{-2\sqrt{2}-\sqrt{2}}{1+\sqrt{2}\cdot \left( -2\sqrt{2} \right)} \right|=\sqrt{2}\ne \sqrt{3}\)，
这说明直线 \(AB\) 与直线 \(l\) 的夹角不是 \(60^\circ\)，矛盾. 因此 C 错误.
对于 D，当 \(k=2\) 时，直线 \(l\) 的方程为 \(y=2x+2\). 设点 \(A\) 到直线 \(l\) 的距离为 \(d\). 因为 \(B\)，\(C\) 在直线 \(l\) 上，且 \(\triangle ABC\) 为等边三角形，所以 \(d\) 是等边三角形的高. 设边长为 \(s\)，则 \(d=\frac{\sqrt{3}}{2}s\)，故 \(s=\frac{2d}{\sqrt{3}}\). 于是三角形面积 \(S_{\triangle ABC}=\frac{1}{2}sd=\frac{d^2}{\sqrt{3}}\). 所以只需求点 \(A\) 到直线 \(y=2x+2\) 的距离的最小值.
设 \(A(x,y)\) 在抛物线上，则 \(y^2=8x\). 点 \(A\) 到直线 \(2x-y+2=0\) 的距离为 \(d=\frac{\left| 2x-y+2 \right|}{\sqrt{5}}\). 由 \(x=\frac{y^2}{8}\)，得 \(2x-y+2=\frac{y^2}{4}-y+2=\frac{\left( y-2 \right)^2}{4}+1>0\). 因此 \(d=\frac{\frac{y^2}{4}-y+2}{\sqrt{5}}\). 当 \(y=2\) 时，\(d\) 取得最小值 \(d_{\min}=\frac{1}{\sqrt{5}}\). 于是 \(S_{\min}=\frac{d_{\min}^2}{\sqrt{3}}=\frac{1}{5\sqrt{3}}\). 所以 \(S_{\min}=\frac{\sqrt{3}}{15}\).
所以 D 正确.
故选：ABD.
\end{solution}
\section{填空题}

\begin{problem}
记 \(S_n\) 为等差数列 \(\{a_n\}\) 的前 \(n\) 项和. 若 \(a_1=-1\)，\(a_3=5\)，则 \(S_6=\fillinblank{}\).

\end{problem}

\begin{answer}
39
\end{answer}

\begin{solution}
设等差数列的公差为 \(d\). 由 \(a_3=a_1+2d\) 得 \(5=-1+2d\)，所以 \(d=3\). 于是 \(S_6=\frac{6}{2}\left[ 2\times \left( -1 \right)+5\times 3 \right]=3\times 13\). 所以 \(S_6=39\).
故填：39.
\end{solution}

\begin{problem}
若函数 \(f\left( x \right)=2^x+2^{2-x}-m\) 有两个零点，则 \(m\) 的取值范围是\fillinblank{}.

\end{problem}

\begin{answer}
\(\left( 4,+\infty \right)\)
\end{answer}

\begin{solution}
令 \(t=2^x\)，则 \(t>0\)，且 \(2^{2-x}=\frac{4}{t}\). 方程 \(f(x)=0\) 化为 \(t+\frac{4}{t}=m\)（\(t>0\)）. 由基本不等式，得 \(t+\frac{4}{t}\ge 2\sqrt{t\cdot \frac{4}{t}}=4\). 当且仅当 \(t=2\) 时取等号. 因此 \(t+\frac{4}{t}\) 在 \(t>0\) 上的值域为 \(\left[ 4,+\infty \right)\).
要使函数 \(f(x)\) 有两个零点，就要使方程 \(t+\frac{4}{t}=m\) 有两个不同的正实根. 这当且仅当 \(m>4\). 故填：\(\left( 4,+\infty \right)\).
\end{solution}

\begin{problem}
已知球 \(O\) 的体积为 \(4\sqrt{3}\pi\)，\(A\)，\(B\)，\(C\)，\(D\) 四点均在球 \(O\) 的球面上，\(\triangle ABC\) 为等边三角形，\(DA=DB=DC=\sqrt{2}\)，则 \(\triangle ABC\) 的面积为\fillinblank{}.

\end{problem}

\begin{answer}
\(\frac{5\sqrt{3}}{4}\)
\end{answer}

\begin{solution}
设球 \(O\) 的半径为 \(R\). 由 \(\frac{4}{3}\pi R^3=4\sqrt{3}\pi\) 得 \(R^3=3\sqrt{3}\)，所以 \(R=\sqrt{3}\).
不妨设球心为坐标原点，且 \(D\left( 0,0,\sqrt{3} \right)\). 设球面上任一点 \(X(x,y,z)\) 满足 \(DX=\sqrt{2}\). 因为 \(X\) 在球面上，所以 \(x^2+y^2+z^2=3\). 又 \(DX^2=x^2+y^2+\left( z-\sqrt{3} \right)^2=2\). 两式相减，得 \(3-2\sqrt{3}z+3=2\)，所以 \(z=\frac{2}{\sqrt{3}}\). 这说明 \(A\)，\(B\)，\(C\) 都在平面 \(z=\frac{2}{\sqrt{3}}\) 上. 因此 \(\triangle ABC\) 的外接圆是该平面与球面的交圆.
该交圆的圆心为 \(\left( 0,0,\frac{2}{\sqrt{3}} \right)\)，半径为 \(r=\sqrt{3-\left( \frac{2}{\sqrt{3}} \right)^2}=\sqrt{\frac{5}{3}}\). 因为 \(\triangle ABC\) 为等边三角形，所以它内接于这个半径为 \(r\) 的圆. 设等边三角形边长为 \(a\)，则 \(r=\frac{a}{\sqrt{3}}\)，故 \(a=\sqrt{3}r=\sqrt{5}\). 于是 \(S_{\triangle ABC}=\frac{\sqrt{3}}{4}a^2=\frac{\sqrt{3}}{4}\times 5\). 所以 \(S_{\triangle ABC}=\frac{5\sqrt{3}}{4}\).
故填：\(\frac{5\sqrt{3}}{4}\).
\end{solution}
\section{解答题}

\begin{problem}
从某厂生产的某种电子产品中随机抽取了若干件进行试验，测试它们首次出现故障的时间（单位：天），由试验结果得到如下频率分布直方图：

\begin{enumerate}
\item 估计这种电子产品首次出现故障的时间的第一四分位数及中位数（假设数据在组内均匀分布）；
\item 记 \(\hat{p}\) 为 \(1\) 件这种电子产品首次出现故障的时间小于 \(365\) 天的概率的估计值.
    \begin{enumerate}
    \item 求 \(\hat{p}\)；
    \item 该厂向某用户销售 \(100\) 件这种电子产品，记 \(X\) 为这 \(100\) 件电子产品中首次出现故障的时间小于 \(365\) 天的件数，假设 \(X\sim B\left( 100,\hat{p} \right)\)，求 \(E\left( X \right)\)，\(D\left( X \right)\).
    \end{enumerate}
\end{enumerate}

\texfigure[max width=8.8cm,max totalheight=5.8cm]{img_repaint/2026/national_paper_2/q15_fig1.tex}

\end{problem}

\begin{solution}
（1）由直方图可知，各组组距都为 \(10\). 各组的频率分别为 \(0.05\)，\(0.10\)，\(0.20\)，\(0.25\)，\(0.15\)，\(0.15\)，\(0.05\)，\(0.05\).
所以累计频率依次为 \(0.05\)，\(0.15\)，\(0.35\)，\(0.60\)，\(0.75\)，\(0.90\)，\(0.95\)，\(1\).
第一四分位数对应累计频率 \(0.25\). 因为 \(0.15<0.25<0.35\)，所以第一四分位数落在区间 \([365,375)\) 内. 在该组内还需要补上频率 \(0.25-0.15=0.10\). 该组的频率密度为 \(0.020\)，所以需要的组内长度为 \(\frac{0.10}{0.020}=5\). 故第一四分位数约为 \(365+5=370\).
中位数对应累计频率 \(0.50\). 因为 \(0.35<0.50<0.60\)，所以中位数落在区间 \([375,385)\) 内. 在该组内还需要补上频率 \(0.50-0.35=0.15\). 该组的频率密度为 \(0.025\)，所以需要的组内长度为 \(\frac{0.15}{0.025}=6\). 故中位数约为 \(375+6=381\).
因此，第一四分位数约为 \(370\) 天，中位数约为 \(381\) 天.

（2）（i）首次出现故障的时间小于 \(365\) 天的频率为前两组频率之和，所以 \(\hat{p}=0.05+0.10=0.15\).

（2）因为 \(X\sim B\left( 100,0.15 \right)\)，所以 \(E(X)=100\times 0.15=15\)，\(D(X)=100\times 0.15\times 0.85=\frac{51}{4}\). 因此 \(E(X)=15\)，\(D(X)=\frac{51}{4}\).
\end{solution}

\begin{problem}
如图，三棱锥 \(A\)-\(BCD\) 中，点 \(E\) 在 \(BD\) 上，\(AE\perp CE\)，\(AE\perp DE\)，\(CD\perp AD\).

\begin{enumerate}
\item 证明：\(CD\perp AB\)；
\item 若 \(DE=2\)，\(BE=1\)，\(AE=\sqrt{2}\)，\(CD=2\sqrt{3}\)，求直线 \(AD\) 与平面 \(ABC\) 所成角的正弦值.
\end{enumerate}

\bitmapfigure[max width=5.6cm,max totalheight=7.0cm]{img/2026/national_paper_2/q16_fig1.png}

\end{problem}

\begin{solution}
（1）因为 \(AE\perp CE\)，\(AE\perp DE\)，而 \(CE\) 与 \(DE\) 是平面 \(CED\) 内两条相交直线，所以 \(AE\perp \text{平面} CED\). 于是 \(AE\perp CD\).
又已知 \(CD\perp AD\). 因为 \(E\) 在 \(BD\) 上，所以 \(B\)，\(D\)，\(E\) 三点共线，从而平面 \(AEB\) 与平面 \(AED\) 是同一个平面. 直线 \(CD\) 同时垂直于平面 \(AED\) 内的两条相交直线 \(AE\) 与 \(AD\)，故 \(CD\perp \text{平面} AED=\text{平面} AEB\). 而 \(AB\subset \text{平面} AEB\)，所以 \(CD\perp AB\).

（2）以 \(E\) 为坐标原点，所在平面 \(CED\) 为 \(xOy\) 平面，建立空间直角坐标系. 由题意可取 \(E(0,0,0)\)，\(D(2,0,0)\)，\(B(-1,0,0)\)，\(A(0,0,\sqrt{2})\). 设 \(C(x,y,0)\).
由 \(CD\perp AD\)，得 \(\overrightarrow{CD}\cdot \overrightarrow{AD}=0\). 其中 \(\overrightarrow{CD}=(x-2,y,0)\)，\(\overrightarrow{AD}=(2,0,-\sqrt{2})\)，所以 \(2\left( x-2 \right)=0\)，即 \(x=2\). 又 \(CD=2\sqrt{3}\)，故 \(\sqrt{\left( 2-2 \right)^2+y^2}=2\sqrt{3}\)，从而 \(\lvert y \rvert=2\sqrt{3}\). 不妨取 \(C(2,2\sqrt{3},0)\).
于是 \(\overrightarrow{AB}=(-1,0,-\sqrt{2})\)，\(\overrightarrow{AC}=(2,2\sqrt{3},-\sqrt{2})\)，\(\overrightarrow{AD}=(2,0,-\sqrt{2})\). 取平面 \(ABC\) 的一个法向量 \(\bs{n}=\overrightarrow{AB}\times \overrightarrow{AC}=\left( 2\sqrt{6},-3\sqrt{2},-2\sqrt{3} \right)\). 设直线 \(AD\) 与平面 \(ABC\) 所成角为 \(\theta\)，则 \(\sin \theta=\frac{\left| \overrightarrow{AD}\cdot \bs{n} \right|}{\lvert \overrightarrow{AD} \rvert \lvert \bs{n} \rvert}\). 又 \(\overrightarrow{AD}\cdot \bs{n}=2\cdot 2\sqrt{6}+0+\left( -\sqrt{2} \right)\left( -2\sqrt{3} \right)=6\sqrt{6}\)，\(\lvert \overrightarrow{AD} \rvert=\sqrt{2^2+\left( \sqrt{2} \right)^2}=\sqrt{6}\)，\(\lvert \bs{n} \rvert=\sqrt{\left( 2\sqrt{6} \right)^2+\left( -3\sqrt{2} \right)^2+\left( -2\sqrt{3} \right)^2}=3\sqrt{6}\). 因此 \(\sin \theta=\frac{6\sqrt{6}}{\sqrt{6}\cdot 3\sqrt{6}}=\frac{\sqrt{6}}{3}\). 所以直线 \(AD\) 与平面 \(ABC\) 所成角的正弦值为 \(\frac{\sqrt{6}}{3}\).
\end{solution}

\begin{problem}
在 \(\triangle ABC\) 中，已知 \(\cos B=\frac{3}{4}\)，\(\cos^2\left( A+C \right)+\sin A\sin C=1\).

\begin{enumerate}
\item 证明：\(\triangle ABC\) 为钝角三角形；
\item 若 \(\triangle ABC\) 的面积为 \(\frac{\sqrt{7}}{4}\)，求 \(\triangle ABC\) 的周长.
\end{enumerate}

\end{problem}

\begin{solution}
（1）因为 \(A+C=\pi-B\)，所以 \(\cos^2\left( A+C \right)=\cos^2 B=\left( \frac{3}{4} \right)^2=\frac{9}{16}\). 由题设得 \(\sin A\sin C=1-\frac{9}{16}=\frac{7}{16}\).
又 \(\sin A\sin C=\frac{\cos\left( A-C \right)-\cos\left( A+C \right)}{2}\). 而 \(\cos\left( A+C \right)=\cos\left( \pi-B \right)=-\cos B=-\frac{3}{4}\)，故 \(\frac{\cos\left( A-C \right)+\frac{3}{4}}{2}=\frac{7}{16}\). 化简得 \(\cos\left( A-C \right)=\frac{1}{8}\).
因为 \(\cos\left| A-C \right|=\frac{1}{8}<\frac{3}{4}=\cos B\)，且 \(0<B<\pi\)，\(0\le \left| A-C \right|<\pi\)，所以 \(\left| A-C \right|>B\). 不妨设 \(A\ge C\)，则 \(A-C>B\). 于是 \(A=\frac{\left( A+C \right)+\left( A-C \right)}{2}\). 又 \(A>\frac{\left( \pi-B \right)+B}{2}=\frac{\pi}{2}\).
因此 \(\triangle ABC\) 为钝角三角形.

（2）设边 \(a\)，\(b\)，\(c\) 分别是角 \(A\)，\(B\)，\(C\) 的对边，设 \(\lambda=2R\)，其中 \(R\) 为外接圆半径. 由正弦定理，\(a=\lambda \sin A\)，\(b=\lambda \sin B\)，\(c=\lambda \sin C\).
由 \(\cos B=\frac{3}{4}\)，得 \(\sin B=\frac{\sqrt{7}}{4}\). 又由面积公式 \(S_{\triangle ABC}=\frac{1}{2}ac\sin B=\frac{\sqrt{7}}{4}\)，得 \(ac=\frac{2S_{\triangle ABC}}{\sin B}=2\). 另一方面，\(ac=\lambda^2\sin A\sin C=\lambda^2\cdot \frac{7}{16}\). 所以 \(\lambda^2\cdot \frac{7}{16}=2\)，即 \(\lambda=\frac{4\sqrt{14}}{7}\).
再由 \(\sin A+\sin C=2\sin \frac{A+C}{2}\cos \frac{A-C}{2}\)，可得 \(\sin \frac{A+C}{2}=\sin \frac{\pi-B}{2}=\cos \frac{B}{2}\). 由半角公式，\(\cos \frac{B}{2}=\sqrt{\frac{1+\cos B}{2}}=\sqrt{\frac{1+\frac{3}{4}}{2}}=\frac{\sqrt{14}}{4}\). 又 \(\cos \frac{A-C}{2}=\sqrt{\frac{1+\cos\left( A-C \right)}{2}}=\sqrt{\frac{1+\frac{1}{8}}{2}}=\frac{3}{4}\). 所以 \(\sin A+\sin C=2\cdot \frac{\sqrt{14}}{4}\cdot \frac{3}{4}=\frac{3\sqrt{14}}{8}\).
于是三角形周长 \(a+b+c=\lambda\left( \sin A+\sin B+\sin C \right)=\frac{4\sqrt{14}}{7}\left( \frac{3\sqrt{14}}{8}+\frac{\sqrt{7}}{4} \right)\). 所以 \(a+b+c=3+\sqrt{2}\).
所以 \(\triangle ABC\) 的周长为 \(3+\sqrt{2}\).
\end{solution}

\begin{problem}
已知椭圆 \(E:\frac{x^2}{a^2}+y^2=1\)（\(a>1\)），过 \(E\) 的右焦点且与 \(x\) 轴垂直的直线被 \(E\) 截得的线段长为 \(\sqrt{2}\).

\begin{enumerate}
\item 求 \(E\) 的离心率；
\item 设 \(O\) 为坐标原点，给定点 \(G\left( t,0 \right)\)（\(t\ne 0\)）. \(A\left( x_0,y_0 \right)\)（\(y_0\ne 0\)）为 \(E\) 上的点，过 \(A\) 作 \(y\) 轴的垂线，垂足为 \(B\)，直线 \(OA\) 与直线 \(GB\) 的交点为 \(P\)，当 \(A\) 在 \(E\) 上运动时，点 \(P\) 的轨迹为 \(M\).
    \begin{enumerate}
    \item 求 \(M\) 的方程；
    \item 当 \(t\) 为何值时，\(M\) 有对称中心？当 \(M\) 有对称中心时，将 \(M\) 平移后得到曲线 \(M'\)，使得 \(O\) 为 \(M'\) 的对称中心，说明 \(M'\) 是什么曲线.
    \end{enumerate}
\end{enumerate}

\end{problem}

\begin{solution}
（1）设椭圆 \(E\) 的右焦点为 \(F(c,0)\). 因为椭圆方程为 \(\frac{x^2}{a^2}+y^2=1\)，所以其短半轴长为 \(1\)，并且 \(c^2=a^2-1\).
过右焦点且与 \(x\) 轴垂直的直线为 \(x=c\). 把 \(x=c\) 代入椭圆方程，得 \(\frac{c^2}{a^2}+y^2=1\). 由于 \(c^2=a^2-1\)，所以 \(y^2=1-\frac{a^2-1}{a^2}=\frac{1}{a^2}\). 因此该直线被椭圆截得的线段长为 \(\frac{2}{a}\). 由题意 \(\frac{2}{a}=\sqrt{2}\)，故 \(a=\sqrt{2}\). 于是 \(c=\sqrt{a^2-1}=1\). 所以椭圆的离心率为 \(e=\frac{c}{a}=\frac{1}{\sqrt{2}}=\frac{\sqrt{2}}{2}\).

（2）（i）由（1）知椭圆 \(E\) 的方程为 \(\frac{x^2}{2}+y^2=1\). 设 \(P(x,y)\). 因为 \(P\) 在直线 \(OA\) 上，所以存在实数 \(\lambda\)，使得 \(P=\left( \lambda x_0,\lambda y_0 \right)\). 又 \(B=(0,y_0)\)，点 \(P\) 在直线 \(GB\) 上. 由分点关系可得 \(x=t\left( 1-\lambda \right)\)，\(y=\lambda y_0\). 所以 \(\lambda=1-\frac{x}{t}\)，从而 \(x_0=\frac{x}{\lambda}=\frac{tx}{t-x}\)，且 \(y_0=\frac{y}{\lambda}=\frac{ty}{t-x}\). 由于点 \(A(x_0,y_0)\) 在椭圆上，故 \(\frac{x_0^2}{2}+y_0^2=1\). 代入上式，得 \(\frac{1}{2}\left( \frac{tx}{t-x} \right)^2+\left( \frac{ty}{t-x} \right)^2=1\). 整理得
\(
\left( t^2-2 \right)x^2+2t^2y^2+4tx-2t^2=0
\)
这就是轨迹 \(M\) 的方程.

（2）对方程 \(\left( t^2-2 \right)x^2+2t^2y^2+4tx-2t^2=0\) 配方. 当 \(t^2-2=0\)，即 \(t=\pm \sqrt{2}\) 时，方程中 \(x\) 只一次出现，曲线 \(M\) 没有对称中心.

当 \(t^2-2\ne 0\) 时，有
\(
\left( t^2-2 \right)\left( x+\frac{2t}{t^2-2} \right)^2+2t^2y^2=\frac{2t^4}{t^2-2}
\)
由配方后的形式可知，曲线 \(M\) 的对称中心为 \(\left( -\frac{2t}{t^2-2},0 \right)\). 因此，\(M\) 有对称中心的充要条件是 \(t\ne \pm \sqrt{2}\).

当 \(t\ne \pm \sqrt{2}\) 时，把曲线平移，使其中心移到原点. 设 \(X=x+\frac{2t}{t^2-2}\)，\(Y=y\).
则 \(M\) 化为
\(
\left( t^2-2 \right)X^2+2t^2Y^2=\frac{2t^4}{t^2-2}
\)

若 \(t^2>2\)，则上式右边为正，且 \(X^2\)，\(Y^2\) 的系数都为正，所以 \(M'\) 是椭圆.

若 \(0<t^2<2\)，则上式两边同乘以 \(-1\)，得
\(
\left( 2-t^2 \right)X^2-2t^2Y^2=\frac{2t^4}{2-t^2}
\)
所以 \(M'\) 是双曲线.

综上，当 \(t\ne \pm \sqrt{2}\) 时，\(M\) 有对称中心；此时平移后的曲线 \(M'\) 在 \(\lvert t \rvert>\sqrt{2}\) 时是椭圆，在 \(0<\lvert t \rvert<\sqrt{2}\) 时是双曲线.
\end{solution}

\begin{problem}
已知函数 \(f\left( x \right)=x\e^x+ax+b\)，曲线 \(y=f\left( x \right)\) 在点 \(\left( 0,f\left( 0 \right) \right)\) 处的切线方程为 \(y=-2x+1\).

\begin{enumerate}
\item 求 \(a\)，\(b\)；
\item 当 \(x>0\) 时，\(f\left( x+m \right)-f\left( x \right)>m\)，求 \(m\) 的取值范围；
\item 当 \(x>0\) 时，\(f\left( k+x \right)+f\left( k-x \right)>2f\left( k \right)\)，求 \(k\) 的最小值.
\end{enumerate}

\end{problem}

\begin{solution}
（1）由切线方程 \(y=-2x+1\) 可知 \(f(0)=1\)，\(f'(0)=-2\). 而 \(f(0)=b\)，所以 \(b=1\).
又 \(f'(x)=\left( x+1 \right)\e^x+a\)，因此 \(f'(0)=1+a=-2\)，故 \(a=-3\).

（2）由（1）知 \(f(x)=x\e^x-3x+1\). 设 \(g(x)=f\left( x+m \right)-f(x)-m\). 则 \(g(x)=\left( x+m \right)\e^{x+m}-x\e^x-4m=\e^x\left[ \left( \e^m-1 \right)x+m\e^m \right]-4m\).

若 \(m<0\)，则 \(1-\e^m>0\). 取 \(x>\max\left\{1,\sqrt{\frac{-4m}{1-\e^m}}\right\}\)，由 \(\e^x>x\)，得
\(
g(x)<-\left( 1-\e^m \right)x\e^x-4m<-\left( 1-\e^m \right)x^2-4m<0
\)
不满足题意. 因此 \(m\ge 0\).

当 \(m=0\) 时，\(g(x)\equiv 0\)，不满足严格大于号，所以 \(m>0\).

对 \(m>0\)，若 \(m\ge \ln 4\)，则 \(\e^m\ge 4\). 当 \(x>0\) 时，\(\e^x>1\)，且 \(\left( \e^m-1 \right)x+m\e^m>m\e^m\)，所以 \(g(x)>m\e^m-4m\ge 0\)，满足题意.

若 \(0<m<\ln 4\)，记 \(C=\e^m\)，\(d=m\left( 4-C \right)>0\). 取 \(x\) 满足 \(0<x<\min\left\{\frac{1}{2},\frac{d}{C\left( 3+2m \right)}\right\}\). 由导数可得，当 \(0<x<\frac{1}{2}\) 时，\(\e^x<1+2x\)，又 \(x^2<x\)，所以
\(
g(x)<C\left( 1+2x \right)\left( m+x \right)-4m<mC-4m+C\left( 3+2m \right)x<0
\)
不满足题意. 因此 \(m\) 的取值范围为 \(\left[ \ln 4,+\infty \right)\).

（3）设 \(h(x)=f\left( k+x \right)+f\left( k-x \right)-2f(k)\). 代入 \(f(x)=x\e^x-3x+1\)，得 \(h(x)=\left( k+x \right)\e^{k+x}+\left( k-x \right)\e^{k-x}-2k\e^k\). 提取公因子 \(\e^k\)，得 \(h(x)=\e^k\left[ \left( k+x \right)\e^x+\left( k-x \right)\e^{-x}-2k \right]\). 再化简为 \(h(x)=2\e^k\left[ k\left( \cosh x-1 \right)+x\sinh x \right]\).
因为 \(\cosh x-1>0\)（\(x>0\)），所以 \(h(x)>0\) 等价于 \(k+\frac{x\sinh x}{\cosh x-1}>0\). 又 \(\frac{x\sinh x}{\cosh x-1}=x\coth \frac{x}{2}\). 令 \(u=\frac{x}{2}>0\)，则 \(x\coth \frac{x}{2}=2u\coth u\). 由 \(\tanh u<u\)（\(u>0\)）可知 \(u\coth u>1\)，故 \(x\coth \frac{x}{2}>2\). 于是当 \(k\ge -2\) 时，对任意 \(x>0\)，都有 \(k+\frac{x\sinh x}{\cosh x-1}>0\)，从而 \(h(x)>0\).

反过来，若 \(k<-2\)，取 \(x\) 满足 \(0<x<-k-2\). 由 \(\e^x>1+x\)，得 \(x\coth \frac{x}{2}=x+\frac{2x}{\e^x-1}<x+2\). 因而 \(k+x\coth \frac{x}{2}<k+x+2<0\)，这时 \(h(x)<0\)，不满足题意.
因此满足条件的最小值为 \(-2\).
\end{solution}
